\section{Tema 6}
\section*{Respuestas del Cuestionario}

\begin{enumerate}[label=\textbf{Pregunta \arabic*:}, leftmargin=*, itemsep=4ex]

\item \textbf{La Balanza de Rentas Secundarias mantenía superávit en el pasado gracias a las remesas enviadas por emigrantes españoles.}
    \begin{itemize}
        \item \textbf{Respuesta:} \textcolor{blue!80!black}{Verdadero.}
        \item \textbf{Justificación:} Hasta 1985 la balanza de rentas secundarias presentaba superávit debido a los envíos (remesas) de los emigrantes españoles.
    \end{itemize}

\item \textbf{La fase expansiva entre 1995 y 2007 se caracterizó por una reducción drástica del gasto de las familias en bienes de consumo duradero.}
    \begin{itemize}
        \item \textbf{Respuesta:} \textcolor{red!80!black}{Falso.}
        \item \textbf{Justificación:} Ocurrió lo contrario. La facilidad en la captación de capital y su bajo coste de financiación provocó un incremento del gasto de las familias en viviendas y bienes de consumo duradero.
    \end{itemize}

\item \textbf{La dotación de capital humano y de infraestructuras han sido factores determinantes para atraer inversión extranjera directa a España.}
    \begin{itemize}
        \item \textbf{Respuesta:} \textcolor{blue!80!black}{Verdadero.}
        \item \textbf{Justificación:} España atrajo inversión extranjera no solo por la adhesión a la UE, sino por contar con un contexto de crecimiento económico, buena dotación de infraestructuras y buena dotación de capital humano.
    \end{itemize}

\item \textbf{La mayor desventaja comercial de España se localiza en los sectores de alimentos y bienes de consumo duradero.}
    \begin{itemize}
        \item \textbf{Respuesta:} \textcolor{red!80!black}{Falso.}
        \item \textbf{Justificación:} España presenta precisamente una ventaja comercial en productos derivados de la agricultura (alimentos) y el automóvil. Su desventaja se encuentra en bienes intermedios energéticos y en bienes de equipo.
    \end{itemize}

\item \textbf{El saldo de la Balanza de Rentas Primarias suele ser positivo debido a que España es un gran emisor de capitales.}
    \begin{itemize}
        \item \textbf{Respuesta:} \textcolor{red!80!black}{Falso.}
        \item \textbf{Justificación:} El saldo ha sido mayoritariamente negativo ya que España es, sobre todo, un país receptor de capitales extranjeros y debe pagar los dividendos e intereses correspondientes al exterior.
    \end{itemize}

\item \textbf{La adhesión a la Unión Europea supuso un estímulo para la llegada de inversiones extranjeras en manufacturas de intensidad tecnológica media-alta y alta.}
    \begin{itemize}
        \item \textbf{Respuesta:} \textcolor{blue!80!black}{Verdadero.}
        \item \textbf{Justificación:} La entrada en la UE reforzó el interés de los países miembros por invertir en el país, dirigiéndose especialmente a este tipo de manufacturas dependientes de importaciones.
    \end{itemize}

\item \textbf{La inversión extranjera en el sector de la energía en España se ha visto impulsada recientemente por proyectos en el ámbito de las energías renovables.}
    \begin{itemize}
        \item \textbf{Respuesta:} \textcolor{blue!80!black}{Verdadero.}
        \item \textbf{Justificación:} Se ha registrado un avance muy notable en la actividad inversora dentro de las energías renovables, ocupando España posiciones muy destacadas globalmente.
    \end{itemize}

\item \textbf{La implementación de los fondos Next Generation EU tras la pandemia supuso una reducción del saldo positivo de la Balanza de Capital debido al aumento de los pagos.}
    \begin{itemize}
        \item \textbf{Respuesta:} \textcolor{red!80!black}{Falso.}
        \item \textbf{Justificación:} La recepción de estos fondos ha provocado un efecto positivo, permitiendo que el saldo de la balanza de capital volviera a situarse en el 1\% del PIB desde 2022.
    \end{itemize}

\item \textbf{Los bienes de equipo representan una de las principales partidas tanto en las exportaciones como en las importaciones españolas.}
    \begin{itemize}
        \item \textbf{Respuesta:} \textcolor{blue!80!black}{Verdadero.}
        \item \textbf{Justificación:} Los bienes de equipo alcanzan cuotas máximas en importaciones (22,4\% en 2023) y también conforman una porción vital en las exportaciones (19,5\%).
    \end{itemize}

\item \textbf{España presenta una ventaja comercial clara en sectores como los productos derivados de la agricultura y el automóvil.}
    \begin{itemize}
        \item \textbf{Respuesta:} \textcolor{blue!80!black}{Verdadero.}
        \item \textbf{Justificación:} El perfil de la ventaja comercial española se sitúa, efectivamente, en los derivados agrícolas (alimentos) y en el sector automovilístico.
    \end{itemize}

\item \textbf{El proceso de devaluación interna tras la crisis de 2008 permitió una reducción de los costes laborales y una mejora de la competitividad-precios.}
    \begin{itemize}
        \item \textbf{Respuesta:} \textcolor{blue!80!black}{Verdadero.}
        \item \textbf{Justificación:} La ''devaluación interna'' redujo costes laborales y la inflación diferencial, lo que incrementó la competitividad vía precios y mejoró el saldo comercial.
    \end{itemize}

\item \textbf{La inversión en cartera se diferencia de la directa en que la primera busca explícitamente el control y la gestión de la sociedad participada.}
    \begin{itemize}
        \item \textbf{Respuesta:} \textcolor{red!80!black}{Falso.}
        \item \textbf{Justificación:} Es al revés. La inversión directa es la que persigue influir significativamente y gestionar la sociedad. La inversión en cartera busca una rentabilidad a corto plazo sin intención de control.
    \end{itemize}

\item \textbf{Un déficit persistente en la balanza por cuenta corriente implica una necesidad de financiación exterior y acumulación de deuda.}
    \begin{itemize}
        \item \textbf{Respuesta:} \textcolor{blue!80!black}{Verdadero.}
        \item \textbf{Justificación:} La teoría señala la necesidad de financiación exterior y la consiguiente acumulación de deuda como las principales implicaciones de un déficit en cuenta corriente.
    \end{itemize}

\item \textbf{La similitud en los niveles de renta entre países es un factor que desincentiva el comercio intraindustrial bilateral.}
    \begin{itemize}
        \item \textbf{Respuesta:} \textcolor{red!80!black}{Falso.}
        \item \textbf{Justificación:} La similitud en el nivel de renta entre países es, de hecho, un factor que fomenta (en sentido positivo) el comercio intraindustrial.
    \end{itemize}

\item \textbf{La tasa de cobertura del comercio de bienes ha superado el cien por cien de forma estructural desde la crisis de 2008.}
    \begin{itemize}
        \item \textbf{Respuesta:} \textcolor{red!80!black}{Falso.}
        \item \textbf{Justificación:} La tasa de cobertura de los bienes ha superado el 90\% desde 2013 (llegando al 86,9\% en 2022). No supera estructuralmente el 100\% salvo que se le añadan los servicios.
    \end{itemize}

\item \textbf{Antes de la integración en la UE, la devaluación del tipo de cambio era una medida habitual para corregir el déficit exterior.}
    \begin{itemize}
        \item \textbf{Respuesta:} \textcolor{blue!80!black}{Verdadero.}
        \item \textbf{Justificación:} Históricamente, cuando el déficit exterior se acercaba al 4\% del PIB, los gobiernos españoles recurrían a la devaluación de la moneda para aumentar la competitividad de las exportaciones.
    \end{itemize}

\item \textbf{La estructura sectorial de las exportaciones españolas se ha transformado radicalmente en la última década, abandonando los bienes intermedios.}
    \begin{itemize}
        \item \textbf{Respuesta:} \textcolor{red!80!black}{Falso.}
        \item \textbf{Justificación:} La estructura sectorial se considera ''estable en el tiempo'' y los bienes intermedios continúan teniendo el mayor peso dentro de las exportaciones.
    \end{itemize}

\item \textbf{Estados Unidos, Reino Unido, Francia y Alemania son los países de origen de más de la mitad del stock de inversión extranjera en España.}
    \begin{itemize}
        \item \textbf{Respuesta:} \textcolor{blue!80!black}{Verdadero.}
        \item \textbf{Justificación:} Estos cuatro países en conjunto representan más del 51\% del stock total de la Inversión Extranjera Directa (IED) recibida en 2023.
    \end{itemize}

\item \textbf{El comercio intraindustrial en España es especialmente relevante en los intercambios con los países de la Unión Europea.}
    \begin{itemize}
        \item \textbf{Respuesta:} \textcolor{blue!80!black}{Verdadero.}
        \item \textbf{Justificación:} El índice de comercio intraindustrial en España subió al 70\% a finales de los años 90, concentrándose sobre todo en los países pertenecientes a la UE.
    \end{itemize}

\item \textbf{La estrategia de futuro de la economía española consiste en concentrar sus flujos comerciales en el mercado de la UE para minimizar riesgos.}
    \begin{itemize}
        \item \textbf{Respuesta:} \textcolor{red!80!black}{Falso.}
        \item \textbf{Justificación:} La estrategia a futuro se basa en la ''diversificación de los mercados'' y no en la concentración. El objetivo es expandirse a economías más dinámicas en Asia, África u Oriente Medio para reducir riesgos.
    \end{itemize}

\item \textbf{La Balanza de Pagos es un documento contable que registra exclusivamente las operaciones comerciales de los residentes de un país con el resto del mundo.}
    \begin{itemize}
        \item \textbf{Respuesta:} \textcolor{red!80!black}{Falso.}
        \item \textbf{Justificación:} Registra tanto las operaciones comerciales como las ''financieras'' de los residentes de un país con el resto del mundo.
    \end{itemize}

\item \textbf{Durante la pandemia de Covid-19, las exportaciones españolas mantuvieron su crecimiento mientras las importaciones caían abruptamente.}
    \begin{itemize}
        \item \textbf{Respuesta:} \textcolor{red!80!black}{Falso.}
        \item \textbf{Justificación:} La pandemia causó que cayeran abruptamente tanto las exportaciones como las importaciones simultáneamente.
    \end{itemize}

\item \textbf{El sector servicios, incluyendo actividades inmobiliarias y banca, lidera la estructura sectorial de la inversión extranjera recibida en España.}
    \begin{itemize}
        \item \textbf{Respuesta:} \textcolor{blue!80!black}{Verdadero.}
        \item \textbf{Justificación:} El sector servicios supone el 49,2\% de la inversión extranjera recibida, con especial protagonismo de las actividades inmobiliarias, comercio y banca/seguros.
    \end{itemize}

\item \textbf{La elasticidad renta de las importaciones españolas es significativamente superior a la elasticidad renta de las exportaciones respecto a la demanda extranjera.}
    \begin{itemize}
        \item \textbf{Respuesta:} \textcolor{blue!80!black}{Verdadero.}
        \item \textbf{Justificación:} Las importaciones dependen de la renta española con una elasticidad igual a 2, mientras que las exportaciones tienen una elasticidad renta respecto a países clientes próxima a 1.
    \end{itemize}

\item \textbf{El esfuerzo inversor de España en el exterior (stock de IED emitida sobre PIB) es marginal en comparación con la media de la Unión Europea.}
    \begin{itemize}
        \item \textbf{Respuesta:} \textcolor{red!80!black}{Falso.}
        \item \textbf{Justificación:} El esfuerzo inversor no es marginal; representa el 39,2\% del PIB (2023), siendo plenamente comparable con los grandes países de la UE y cercano a la media comunitaria.
    \end{itemize}

\end{enumerate}
